\subsection{A fixed-prime finite-field separation premise fails}

The counterexample in \citet[Remark~10(i), pp.~742--743]{LitHoshi2009Configurations},
attributed there to Tamagawa, concerns distinct interior points on a proper
hyperbolic curve.  We give the corresponding construction for the tripod,
with its actual based action, and distinguish geometric conjugacy from
conjugacy in the whole arithmetic group.

\begin{result}{E.150}{Proved}{Distinct finite-field points with the same full pro-\(3\) section}
Let \(U=\mathbb P^1\setminus\{0,1,\infty\}\) over
\(K=\mathbb F_q\), where \(3\nmid q\), and let
\(F=\pi_1(U_{\overline K})^{(3)}\).  Set
\(e=\operatorname{ord}_{\mathbb F_3^\times}(q)\in\{1,2\}\).
If \(L=\mathbb F_{q^d}\) and \(\gcd(d,3e)=1\), then the
full geometrically pro-\(3\) sections associated to any two points in
one \(\operatorname{Gal}(L/K)\)-orbit in \(U(L)\) are conjugate
by an element of \(F\).  The geometric group \(F\) is fixed;
only the arithmetic base field is extended.

In particular, both of the following are explicit failures of injectivity:
\begin{enumerate}
\item Let \(K=\mathbb F_{25}\) and
\(L=\mathbb F_5[\alpha]/(\alpha^4+2)=\mathbb F_{625}\).
The distinct points \(\alpha,-\alpha\in U(L)\) have the same
full pro-\(3\) section class.
\item Let \(K=\mathbb F_5\) and
\(L=\mathbb F_5[\beta]/(\beta^5-\beta-1)=\mathbb F_{3125}\).
The five distinct points \(\beta+i\), \(i\in\mathbb F_5\),
have the same full pro-\(3\) section class.  Here
\(\#L\equiv-1\pmod3\).
\end{enumerate}
These examples disprove a general fixed-pro-\(3\) interior-point
separation assertion, including one restricted to finite fields of
cardinality \(-1\pmod3\).  They do not identify an interior point
section with a tangential section, give an example over a prime field,
or settle the gate in E.135 or the ordinary completed-pair question over
\(\mathbb Q\).
\end{result}
\begin{proof}
Use the unit tangential basepoint at \(0\) to split the arithmetic
extension.  With \(G_K=\widehat{\mathbb Z}\) generated by arithmetic
Frobenius \(\sigma\), this gives an actual continuous action
\(\rho:G_K\to\operatorname{Aut}(F)\) and
\[
 \Pi_K=F\rtimes_\rho G_K,
 \qquad
 \Pi_L=F\rtimes_\rho G_L\subseteq\Pi_K,
 \qquad G_L=d\widehat{\mathbb Z}.
 \tag{E.150.1}
\]
The Frattini quotient is
\(F/\Phi(F)\simeq\mathbb F_3(1)^2\); hence the image of
\(\sigma\) on it is the scalar \(q\), of order \(e\).
The kernel of
\(\operatorname{Aut}(F)\to\operatorname{Aut}(F/\Phi(F))\)
is pro-\(3\).  Indeed, for each finite characteristic \(3\)-group
quotient \(P\) of \(F\), its image lies in the kernel of
\(\operatorname{Aut}(P)\to\operatorname{Aut}(P/\Phi(P))\),
which is a \(3\)-group by the finite Burnside basis theorem; these
quotients detect the topology of \(\operatorname{Aut}(F)\).

Consequently the procyclic group
\(\Gamma=\rho(G_K)\) has a pro-\(3\) subgroup of index \(e\).
The condition \(\gcd(d,3e)=1\) makes the \(d\)-th power map
an automorphism of \(\Gamma\).  Thus
\[
 \rho(G_L)=\rho(G_K),
 \qquad
 \text{there exists }a\in G_L\text{ with }
 \rho(a)=\rho(\sigma).
 \tag{E.150.2}
\]
Write \(\tau:G_K\to\Pi_K\) for the tangential splitting.
Since \(G_K\) is abelian, \(\tau(\ker\rho)\) is central in
\(\Pi_K\).  It follows that conjugation by \(\tau(a)\) and
by \(\tau(\sigma)\) agree on \(\Pi_K\).

For \(x\in U(L)\), functoriality of point decomposition groups
identifies conjugation by \(\tau(\sigma)\) with the action
\(x\mapsto\sigma x\), up to geometric conjugacy.  By changing the
representative of the section of \(\sigma x\), we may therefore write
\[
 s_{\sigma x}(g)=\tau(a)s_x(g)\tau(a)^{-1}
 \quad(g\in G_L).
\]
This is initially a conjugacy inside \(\Pi_L\), not yet the
equivalence relation used for sections.  The element
\[
 h=\tau(a)s_x(a)^{-1}\in F
\]
repairs that distinction: because the image \(s_x(G_L)\) is abelian,
\[
 h s_x(g) h^{-1}
 =\tau(a)s_x(a)^{-1}s_x(g)s_x(a)\tau(a)^{-1}
 =s_{\sigma x}(g).
 \tag{E.150.3}
\]
Iteration proves the orbit assertion in the full group, and hence
simultaneously in all its finite and nilpotent quotients.

For the first example, \(T^4+2\) has no root in \(\mathbb F_5\):
every nonzero fourth power is \(1\).  A quadratic factorization would
have the form
\[
 (T^2+aT+b)(T^2-aT+c),
 \quad bc=2,\quad b+c=a^2,\quad a(c-b)=0.
\]
If \(a=0\), then \(c=-b\) and \(b^2=3\), impossible in
\(\mathbb F_5\); if \(a\ne0\), then \(c=b\) and \(b^2=2\),
also impossible.  This proves irreducibility.  In the resulting field,
\[
 \alpha^{25}=\alpha(\alpha^4)^6
 =3^6\alpha=-\alpha.
 \tag{E.150.4}
\]
Thus \(\alpha\notin\mathbb F_{25}\), its orbit over
\(\mathbb F_{25}\) has size two, and both orbit members lie in
\(U(L)\).  Here \(e=1\) and \(d=2\).

For the second example, any root \(\beta\) of
\(T^5-T-1\) satisfies
\(\beta^{5^j}=\beta+j\) for \(j=0,\ldots,5\).
The least positive \(j\) fixing \(\beta\) is therefore \(5\).
This proves that its minimal polynomial has degree five, hence that
\(T^5-T-1\) is irreducible.  Its five roots are \(\beta+i\),
none in \(\mathbb F_5\), so none equals \(0\) or \(1\).
Here \(e=2\), \(d=5\), and \(5^5\equiv-1\pmod3\).

Finally, the orbit construction cannot change an interior point into
a cusp: the boundary \(\{0,1,\infty\}\) is stable under the
arithmetic action.  This is the precise reason the counterexamples do
not resolve the interior-versus-tangent comparison in E.135.
\end{proof}

\paragraph{Relation to the published counterexample.}
Hoshi's Remark~10(i) first makes the outer action of a proper curve
pro-\(\ell\) by finite base extension and then uses the central
prime-to-\(\ell\) arithmetic factor.  The proof above uses the
tripod's specified tangential splitting, checks the actual action, and
allows its residual scalar of order two.  Thus it supplies a direct
affine example, rather than assuming that the proper-curve example
survives puncturing.  Frobenius-preserving reconstruction of configuration
spaces in \citet[Theorem~0.1, pp.~662--663; Proposition~2.5(ii),
pp.~696--697]{LitHoshi2009Configurations} is compatible with the failure
of injectivity: preservation of the set of point-section classes does
not assert that different points determine different classes.

\paragraph{The stronger containment statement has different hypotheses.}
For comparison, \citet[Theorem~1.3(iv), author p.~6]{LitMochizuki2005Sections}
does exclude an interior decomposition group being contained in a cusp
decomposition group.  Its standing fields are finite extensions of
\(\mathbb Q_p\) (\S0, author p.~2; opening of \S1, p.~4), and its
groups are full profinite.  The proof passes to a cover with hyperbolic
proper compactification and uses point determination there.  These
hypotheses do not yield the singleton-prime finite-field assertion
needed in E.135.  In the finite-field setting,
\citet[Proposition~1.8(ii), pp.~144--145]{AuditSaidiTamagawa2009}
distinguish whole cusp and interior decomposition groups up to
commensurability by their different geometric inertia.  A complement
\(\widehat{\mathbb Z}\subset\mathbb Z_3(1)\rtimes\widehat{\mathbb Z}\)
is not commensurate with the containing group, so that conclusion does
not exclude the relevant containment.
